\documentclass[a4paper,12pt,Times]{article}
\usepackage{abakos} 
\usepackage{booktabs}
\usepackage{bigstrut}
\usepackage{tabls}
% Configurando layout para mostrar codigos Java
\usepackage{listings}
\lstset{
	language=Java,
	 basicstyle=\ttfamily\scriptsize, 
	keywordstyle=\color{purple}\bfseries,
	stringstyle=\color{red},
	commentstyle=\color{green},
	morecomment=[s][\color{blue}]{/**}{*/},
	extendedchars=true,
	showspaces=false,
	showstringspaces=false,
	numbers=left,
	numberstyle=\tiny,
	breaklines=true,
	backgroundcolor=\color{cyan!10},
	breakautoindent=true,
	captionpos=b,
	xleftmargin=18pt,
	xrightmargin=18pt,
	tabsize=4
}
% Fim Java

\newcommand{\monog}{Refatoração em um Sistema Legado}
\newcommand{\tipo}{Artigo}  
\newcommand{\origem}{Brasil}
\newcommand{\editorial}{}  

\newcommand{\AutorA}{Julio Cesar do Carmo Heredia}
\newcommand{\funcaoA}{Aluno}
\newcommand{\emailA}{jul.heredia@gmail.com}
\newcommand{\cursA}{Trabalho de Conclusão de Curso em Sitemas de Informação da PUC Minas}

\newcommand{\AutorB}{Hugo Bastos de Paula} 
\newcommand{\funcaoB}{Orientador}
\newcommand{\emailB}{hugo@pucminas.br}
\newcommand{\cursB}{Instituto de Ciências Exatas e de Informática da PUC Minas}

\newcommand{\univ}{Pontifícia Universidade Católica de Minas Gerais}

\newcommand{\keyword}[1]{\textsf{#1}}


\begin{document}

	\begin{flushleft}
		\vspace{0cm} {
			\singlespacing \Large{\monog \\ }
		}
	\end{flushleft}
	
	\begin{flushright}
		\singlespacing 
		\normalsize{\AutorA \footnote{\funcaoA, \cursA, \origem -- \emailA }} \\
		\normalsize{\AutorB \footnote{\funcaoB, \cursB, \origem -- \emailB }} \\
		%\normalsize{\AutorB \footnote{\funcaoB, E-mail:\emailB \\ \cursB, \origem. }} \\


	\end{flushright}
	\thispagestyle{empty}
	
	\begin{abstract}
		\noindent
			Este trabalho aborda a refatoração de código legado em um sistema de venda de passagens escrito em Java. Esta abordagem se apoia na dificuldade dos desenvolvedores ao encontrar um código em que a sua estrutura original esteja com a sua integridade comprometida e com tecnologias obsoletas. São apresentadas neste trabalho, técnicas de refatoração que levaram este código a ter principalmente mais inteligibilidade, para que seja um sistema de fácil compreensão para seus desenvolvedores. O trabalho também tem o intuito de quantificar a melhoria do código refatorado. O código será metrificado com o auxílio das ferramentas SonarQube e Metrics. Essas ferramentas farão o comparativo do antes e pós refatorações, com o propósito de relatar a melhoria na qualidade de software após o término das refatorações. Foram avaliados dois projetos, o projeto Cliente e o Servidor. O projeto Cliente teve uma melhor margem na melhoria da qualidade de software em relação ao projeto Servidor, onde ouve uma melhoria significativa nas métricas relacionadas à complexidade do código, de 14,44\% na análise do SonarQube e de 44,26\% na complexidade ciclomática no relatório do Metrics.
		\\\textbf{\keyword{Palavras-chave: }} Refatoração, Técnicas de Refatoração, Sistemas Legados, Qualidade de Software e Métricas de Software.
	\end{abstract}
	
%	\newpage 
%	\selectlanguage{english}
%	\begin{abstract}
%		\noindent
%			This paper discusses the refactoring in legacy code of a ticket sales system written in Java. This approach is based on the difficulty of the developers when they find a code that the original structure is compromised with its integrity and obsolete technologies. It is presented in this paper, refactoring techniques that led this code to be mostly more intelligibility, to make it a user-friendly system to its developers. The work also have the purpose of quantifying the improved Refactored code. The code will metrical with the help of some tools. These tools will make the comparison of before and after refactoring, In order to report the improvement in the quality of software after the end of the refactorings.
%		\\\textbf{\keyword{Keywords: }} Refactoring, Refactoring techniques, Legacy Systems, Software Quality and Software Metrics.
%	\end{abstract}
	
	\selectlanguage{brazilian}
	\onehalfspace  % espaçamento 1.5 entre linhas
	\setlength{\parindent}{1.25cm}
	
	\include{textos}
	
	\newpage
	\singlespace{
		\bibliography{bibliografia}
	}
	
\end{document}